\section{续筛新增的数值证据：先比较算法，再比较现象}
\label{sec:remaining-numerics}
本节列出此次全文核读直接影响判断的计算条件。所有物理推导在第一章；这里讨论实际求解与验证，不能把理论写出的方程等同于代码中实际保留的项。

\subsection{全局尺寸扫描应固定什么}
Huang等\citep{F600}用$L_x=25.6\ell,L_z=12.8\ell,\delta=0.5\ell$的Harris片，改变$\ell/d_i$，因此保持的是宏观形状和$\delta/L$，不是固定微观片厚后单独加长盒子。PIC扫描到$\ell/d_i=30$，Hall MHD到100。PIC使用质量比25、20格/$d_i$及每种每格平均250粒子；Hall流体在中心区用$\Delta x=0.1d_i,\Delta z=0.02d_i$，超电阻和黏性采用固定离子尺度归一值。

比较时间也须一致：相同$t\Omega_{ci}$不等于相同全局演化阶段。该文比较相同重联磁通区间并以初始$B_0V_A$归一。但实际核读v1的正文平均区间写$0.05$--$0.3$，图1阴影写约$0.5$--$3$；两者相差十倍。在原始脚本或作者澄清前，可引用已展示的尺寸趋势，不能自称精确复算其平均率。另有一个PIC尺寸的重复运行相差约12\%，应同时记录随机扰动敏感性。

磁岛合并的尺寸扫描\citep{L241}同样改变$\lambda/d_i$并保持形状，但系统由宏观岛供应磁通，与长Harris片并不相同。其率使用$B_m=0.353B_0$与$V_{Am}$归一，平均窗口为$1.5\tau_A$。直接把其约0.6的峰值与其他文章的0.1放在一起，会混淆磁场基准、局地/全局和峰值/平均。

\subsection{非对称层中，改变耗散项可以改变整个宏观演化}
Aunai等\citep{F625}用相同混合模型更换网格控制、电阻和超电阻处理。离子由Boris推进，电子无质量且等温，电磁场二阶差分和预测校正。非对称两侧$|B|=0.5,1.5$，密度$1,1/3$，导引场1；$T_i/T_e=5$。其结果不能简单写成“Hall打开了快重联”：
\begin{enumerate}
 \item 网格控制时，非对称电子层反复伸长并数值撕裂；细化网格改变起岛时刻、层长和平均率。对称算例较不敏感，不能代表所有几何。
 \item 均匀电阻可保持短层，但同时扩散更大尺度的磁能；超电阻更集中在细尺度，其宏观率也不同。
 \item 周期域太短时，返回的流和场可加宽电流层，使磁岛不再形成，或者只出现一次后便维持短X线。看起来稳定的解可能由边界支持。
\end{enumerate}
这是对“任意小尺度断线机制都等效”的反例，不是对某个超电阻闭合已等价于真实动力学的证明。该文部分加密组同时改变了盒长，故严格复现还需分别改变网格与域长，并只比较边界返回影响之前的时间窗。

\subsection{磁岛与分辨率：两类对照必须分开}
Garcia Morillo与Alexakis的正式版本\citep{F788}研究的是二维不可压\emph{电阻MHD}，并没有Hall项。双周期$2\pi$方域、Orszag--Tang初值、$\nu=\eta$、初始磁/动能$0.5/0.125$，另加$k<16$、能量约0.25\%的扰动；GHOST伪谱、四阶Runge--Kutta、$2/3$去混叠，$N=1024$到32768。欠分辨时出现磁岛和近乎平台的率；足够分辨的相应算例中没有这些岛，率更接近Sweet--Parker标度。最高已分辨局地$S_L$约$5.4\times10^5$。作者还讨论了电流片寿命、扰动衰减和被排走的时间，正式版并未证明“所有高$S$磁岛都是假的”。

该文用电流谱在截断波数处相对峰值充分衰减作实用检查，但其经验谱阈值不能直接移植为有限体积或DG的统一收敛判据。本轮已下载并核读10页期刊版本，原8页预印本保留，独立论文仍只算一篇。

Akutagawa等\citep{F606}则比较PIC与电阻MHD。主要PIC质量比1、2、4、8、16、32，空间步长$0.1d_e$，$100d_i\times50d_i$，部分加倍域；主文低质量比易出现次级岛，但这些岛并未像MHD例那样显著提高全局能量转换。附录显示一些粗网格磁岛在加密后消失。真实质量比结论仍是外推：附录虽给质量比100，使用的是$0.5d_e$粗网格，也不等于完成质量比1836验证。

两篇文章分别涉及不同方程、噪声和形成过程，不能合并成一个否定磁岛机制的普遍结论。更有用的共同要求是同时查看磁通拓扑、场谱、网格/时间步、扰动、边界返回、守恒和非理想项预算，而不是仅比一张电流密度彩色图。

\subsection{回旋流体的制造解、超黏性和实际算例是三种检查}
GREENY v1\citep{E473}给出极化和Ampère反演的制造解，展示SOR约$h^2$、其FFT或PCG/Teague组合约$h^4$误差；这里的阶数包含其具体离散，不表示所有谱方法只有四阶。算法使用Arakawa括号和Karniadakis多步，极化求解与磁势求解可选择不同近似。

近边界重联算例中，超黏性太强会抹平初始扰动，太弱会出现格点颗粒结构；部分主峰重合而前期岛形成和起始时间仍变。必须区分：制造解是否正确、守恒误差是否小、特定非线性现象是否随网格及耗散一起收敛。不能只展示其中一项便宣布物理已收敛。后续Full-F研究\citep{L314}的$1024^2$与$32\rho_s$域及迎风/超扩散比较提供了进一步信息，但文中若干参数标注矛盾仍需保留，见第\ref{sec:remaining-audit}节。

\subsection{近期大尺度混合湍流：电子尺度分辨尚需特别核验}
F655\citep{F655}采用$100\pi d_i$方域、$3072^2$空间格、$51^3$离子速度格，质量比100，等温惯性电子，初始注入尺度约$50d_i$。它确实扩大了相对注入尺度，但由公开数值可直接算
\begin{equation*}
 \Delta x={100\pi\over3072}d_i\simeq0.1023d_i,
 \qquad d_e=0.1d_i,\qquad\Delta x/d_e\simeq1.023.
\end{equation*}
电子惯性长度约一格，且使用格点尺度滤波，正文没有独立加密/滤波收敛组。宽度数个$d_i$的涡链可见，并不自动保证断线电子层已收敛。报告因此保留其跨尺度生成机制，同时降低电子层功率峰值和发生率外推的证据等级。

\subsection{RFP计算应把壁面通量写入验证}
NIMROD扩展MHD松弛研究\citep{F798}用$r$--$z$面$120\times64$个双五次谱元、方位$m=0$--5、轴向周期，比较有无Hall及离子回旋黏性；$S=20000$、$d_i/a\simeq0.17$，电子质量放大5倍。它没有求完整热预算，温度方程忽略电阻与黏性加热。

新生成的总动量主要来自壁面的黏性与回旋黏性耦合，再由Maxwell应力向内输运。作者还量化了省略平衡pinch流导致的额外动量损失；特定事件中相对于$r/a=0.7$处Maxwell通量约10\%到22.5\%。因此核心速度变化方向的比较有价值，但不能把总动量增长当作已解释了真实实验的无外源自发加速。
